%Please do not modify this part%
%%%%%%%%%%%%%%%%%%%%%%%%%%%%%%%%

\documentclass[a4paper,10pt]{amsart}  

\textwidth16cm \textheight20.5cm \oddsidemargin-0.1cm
\evensidemargin-0.1cm

\usepackage[utf8]{inputenc}
\usepackage[T1]{fontenc}
\usepackage{amsthm}
\usepackage{amsmath}
\usepackage{amssymb}
\usepackage[inline]{enumitem}
\usepackage[dvips]{graphicx}
\usepackage{color}
\usepackage{comment}
\usepackage{hyperref}
\usepackage{url}
\usepackage{stackrel}
\usepackage{mathrsfs}
\usepackage{stmaryrd}
\usepackage{hyperref}
\usepackage{fancyhdr}
\usepackage[english]{babel}

\usepackage{background} % Per la filigrana

% Configurazione filigrana
\backgroundsetup{
  scale=0.5,            % dimensione logo
  color=black,           % colore
  opacity=0.08,           % trasparenza (0=trasparente, 1=opaco)
  contents={\includegraphics[width=\paperwidth,height=\paperheight,keepaspectratio]{logo.png}},
  position={current page.center} % posizione centrata
}

\definecolor{blue-url}{RGB}{0,0,100}

\hypersetup{
	pdftitle={The Open Seminar},
	pdfmenubar=false,
	pdffitwindow=true,
	pdfstartview=FitH,
	colorlinks=true,
	linkcolor=blue-url,
	citecolor=blue-url,
	urlcolor=blue-url
}

\begin{document}

% The total length (including the references) should not exceed one page.
\thispagestyle{empty}

% 1. Title of the talk
\begin{center}
    \textbf{\Large{Trace-Metrized Algebras}}
\end{center}

\vskip 0.2cm

% 2. Speaker's name
\begin{center}
    Artem~\textsc{Lopatin}
\end{center}

\vskip 0.8cm

% 3. Abstract
\textit{Abstract.} We study finite-dimensional algebras that are not necessarily associative,
commutative, or unital, focusing on symmetric nondegenerate associative
bilinear forms constructed from traces of multiplication operators.
Algebras admitting such forms are called \emph{trace-metrized}.

Our main goal is to determine when an algebra is trace-metrizable,
describe all its trace metrics, and investigate when every metric is a
trace metric. We establish a canonical representation of symmetric
trace forms and show that trace metrizability implies semisimplicity.
We also study how metrizability and trace metrizability behave under
mutations of the multiplication, obtaining criteria involving Lie and
Jordan algebras.

Special attention is given to algebras of dimensions two and three.
We obtain criteria for trace metrizability and describe all trace metrics
in these dimensions. Finally, we construct examples illustrating the
differences between metrizability, trace metrizability, and semisimplicity.

This is joint work with Mar\'{\i}a Alejandra \'{A}lvarez.

\vskip 0.5cm

% 4. Academic bio
\textit{About the speaker.} Artem Lopatin is an Associate Professor of Mathematics at the University
of Campinas (UNICAMP), Brazil. He received his Ph.D. from Omsk State
University in 2004 and his Doctor of Sciences degree from the Sobolev
Institute of Mathematics in 2013.

His research interests include invariant theory, polynomial identities,
representation theory, and nonassociative algebras. His work focuses
on matrix and separating invariants, Lie and Novikov algebras, Weyl
algebras, and octonions. He has held research positions at the Sobolev
Institute of Mathematics and Bielefeld University, as well as a visiting
appointment at the Technical University of Munich. He is an editor of
\emph{Communications in Algebra}. His papers have appeared in journals
including \emph{Journal of Algebra}, \emph{Algebra \& Number Theory},
and \emph{Journal of Pure and Applied Algebra}.

\vskip 0.5cm

% 5. References are optional and are omitted for this seminar abstract.

% 6. Affiliation and address
{\footnotesize
\textsc{Department of Mathematics, University of Campinas
| 13083-859 Campinas, Brazil}}

% 7. Contact details
{\footnotesize
\textit{E-mail}: \texttt{lopatin@unicamp.br}
$\cdot$
\textit{Web}: \url{https://sites.google.com/view/artemlopatin}
}

\end{document}
